% App. G -- extended related work and standards crosswalk (W7)
\section{Extended related work and standards crosswalk}
\label{app:relwork}

\paragraph{VLA policies and how they are evaluated.}
End-to-end policies, from RT-2 \citep{ref3} and OpenVLA \citep{ref4} to \pizero{} \citep{ref62}, FAST \citep{ref64}, \pifive{} \citep{ref63} and GR00T N1 \citep{ref5}, learn from corpora such as DROID \citep{ref65} and are scored on task success \citep{ref6,ref7,ref14,ref66}.
In that frame, sweeping through a bystander and clearing them score the same if both deliver the object.

\paragraph{Trajectory-level safety suites.}
Table~\ref{tabG:related} extends Table~\ref{tab:related}, without the instruction-level SafeAgentBench.
Most concurrent suites score trajectory predicates with no person present: temporal-logic clauses \citep{ref27,ref30}, safety categories across embodiments \citep{ref23}, object constraints and damage \citep{ref69,ref71,ref72}, and hazards for a mobile robot \citep{ref21}.
SafeStage \citep{ref61} is closest in framing, with a during-execution stage on the same simulated DROID arm and \pifive{} among its policies, but its predicates are events on objects and regions.
LIBERO-Safety \citep{ref24} and HRIBench \citep{ref68} place a person, as a perturbable hand proxy or an animated intruder, so a person in the scene, and one who moves, are not our contribution.
Surveys \citep{ref28,ref29}, formal planner-level evaluation \citep{ref31} and guardrail designs \citep{ref82} frame the wider field.

\begin{table}[!h]
\centering
\caption{\textbf{No other suite pairs person-referenced predicates with a blind control and a witness.}
Our reading; ``$\approx$T$n$'': closest sub-type; (\textit{n}): policies. None carries a hazard past a bystander.}
\label{tabG:related}
\scriptsize
\setlength{\tabcolsep}{2.4pt}
\renewcommand{\arraystretch}{0.95}
\begin{tabular}{@{}>{\raggedright\arraybackslash}p{2.35cm} >{\raggedright\arraybackslash}p{1.25cm} >{\raggedright\arraybackslash}p{2.1cm} >{\raggedright\arraybackslash}p{3.0cm} >{\raggedright\arraybackslash}p{1.35cm} >{\raggedright\arraybackslash}p{1.35cm} >{\raggedright\arraybackslash}p{1.3cm}@{}}
\toprule
Suite & Robot & Person in scene & Predicates ($\approx$ our sub-type) & Speed\,/\,force scored & Ablations & Control, witness\\
\midrule
SafeVLA-Bench (9) \citep{ref27} & fixed arm & no & self-contact $\approx$T2, held-object tilt $\approx$T4, contact force & force proxy only & none & no\\
LIBERO-Safety (10) \citep{ref24} & fixed arm & hand proxy, perturbed & collision margin $\approx$T1/T2, perturbed hand $\approx$T6a & no & mitigation demo & no\\
SafeVLA \citep{ref21} & mobile & no & environmental hazards & no & safe-RL method & no\\
SafeManip (6) \citep{ref30} & fixed arm & no & grasp and release $\approx$T4, contact & no & prompt & no\\
ForesightSafety-VLA (4) \citep{ref23} & 5 robots & no & force/torque $\approx$T5b, spatial boundary $\approx$T1 & force only & no & no\\
HazardArena (4) \citep{ref19} & fixed arm & static, unscored & semantic unsafe twins & no & defence & no\\
R2HandoverSim \citep{ref34} & fixed arm & receiver & handover orientation $\approx$T3 & no & no & no\\
SafeStage \citep{ref61} & fixed arm & no & during-execution events $\approx$T1/T2 & no & no & no\\
HRIBench \citep{ref68} & fixed arm & animated & collision-free rate $\approx$T6a & no & no & no\\
VLA-Arena \citep{ref69} & fixed arm & no & hazard avoidance $\approx$T1, state preservation $\approx$T4 & no & no & no\\
MANIGUARD \citep{ref71} & fixed arm & no & spill, topple, drop $\approx$T4 & no & fine-tuning & safe demos\\
OopsieVerse \citep{ref72} & fixed arm & no & object damage & objects & no & no\\
Assistive care \citep{ref70} & fixed arm & patient & region-wise contact force $\approx$T5b & intended & no & no\\
\midrule
\textbf{This work} (5) & Franka; G1 & bystander, hand, passer-by & six sub-types, eight predicates, four dimensions & T5a (G1); tabletop exposure & wording, hidden person & both (tabletop)\\
\bottomrule
\end{tabular}
\end{table}

\paragraph{Handover and human-aware motion.}
Presenting a carried object handle-first is an established handover convention \citep{ref53,ref55,ref25,ref59}, extended to receiver-centred, adaptive and language-grounded orientation \citep{ref32,ref33,ref26} and benchmarked in simulation \citep{ref34}.
These works study a cooperating receiver; T3 asks the question about a bystander who receives nothing.
Human-aware planners build person-referenced motion in \citep{ref49,ref51}; T1, T2 and T6a ask whether a VLA learned it from data, which human-aware demonstrations can supply \citep{ref73}.

\paragraph{Classical machinery and remedies.}
Constraining how a motion unfolds is not new: potential fields \citep{ref8}, control barrier functions \citep{ref9}, shielding \citep{ref10}, safe learning \citep{ref16} and latent safety filters \citep{ref17,ref18} all do it.
The HRI-safety line supplies person-referenced quantities \citep{ref46,ref48,ref47,ref60,ref52}, and Habitat 3.0 reports a human-collision rate \citep{ref45}.
VLA-specific remedies include constraint layers \citep{ref22}, attention-guided and flow-matching filters \citep{ref41,ref42,ref43}, barrier functions for ISO\,10218 separation \citep{ref44}, whole-body filters \citep{ref74,ref75,ref76} and a safe-control kit on the G1 \citep{ref40}.
The benchmark measures how often a policy needs such a remedy.

\paragraph{Standards crosswalk.}
Table~\ref{tabG:standards} maps each sub-type to its governing requirement; where no standard fixes a value, the threshold is ours and its sensitivity is reported (App.~\ref{app:cards}).
ISO/TS\,15066 \citep{ref11}, now in ISO\,10218-1/-2:2025 \citep{ref35,ref36}, governs separation and contact; ISO\,13855 \citep{ref20} names the separation terms and ISO\,13482 \citep{ref37} covers personal-care robots.
Peer suites score object- or hand-referenced versions of T1, T2, T4, T5b and T6a, and two score contact with a person: HRIBench counts collisions with an animated person ($\approx$T6a) and the assistive-care benchmark scores contact force on a patient ($\approx$T5b). None scores T3 toward a bystander, T5a against a person, or T6b.

\begin{table}[!ht]
\centering
\caption{\textbf{Each sub-type maps to a requirement or convention.}
Our reading; ISO/TS\,15066:2016 clauses. SSM: separation monitoring, no contact; PFL: force limiting, blunt contact.}
\label{tabG:standards}
\scriptsize
\setlength{\tabcolsep}{2.6pt}
\renewcommand{\arraystretch}{0.95}
\begin{tabular}{@{}>{\raggedright\arraybackslash}p{2.5cm} >{\raggedright\arraybackslash}p{6.6cm} >{\raggedright\arraybackslash}p{2.9cm} >{\raggedright\arraybackslash}p{1.2cm}@{}}
\toprule
Sub-type & Governing requirement & Measured quantity & Regime\\
\midrule
T1 payload path & SSM (5.5.4); ISO\,13482 hazardous contact & carried hazard to a keep-out & SSM\\
T2 body sweep & PFL for the robot body; the uncertainty term $Z$ of SSM & link to body surface & PFL if blunt\\
T3 hazard presentation & no clause; handover convention \citep{ref25} & hazardous end to the bearing & SSM\\
T4 load tilt & ISO\,13482 hazards from the load and from dropped objects & vessel tilt (spill proxy) & neither\\
T5a speed near a person & SSM (5.5.4) inverted to $v_{\mathrm{allow}}(d)$; ISO\,13855 terms & speed vs separation & SSM\\
T5b contact force & PFL (5.5.5); Annex~A body-region limits & peak force per body region & PFL if blunt\\
T6a crossing hand & SSM with the human-velocity term; protective stop & contact with a moving hand & SSM (stop)\\
T6b passer-by & the slowing that SSM's human-velocity term assumes & speed at closest approach & SSM\\
\bottomrule
\end{tabular}
\end{table}
